%% =================================================================
%% Kingma, Ba. "Adam: A Method for Stochastic Optimization."
%% ICLR 2015. arXiv:1412.6980v9.
%%
%% Reconstruction of page 9 (printed p. 9) as study notes.
%% Generated by: reproduce-all-pages-basic-tex (batch 1-15)
%%
%% Structure: Algorithm 2 box (AdaMax pseudocode) + derivation chain
%% eqs (8)-(12) showing u_t = max(β2·u_{t-1}, |g_t|) + Sec 7.2
%% TEMPORAL AVERAGING + Sec 8 CONCLUSION (starts).
%% Typesetting anomaly preserved: "Initalization" (with missing 'i')
%% in the Sec 7.2 prose — original typo, not fixed per NOTES.md §5.
%% Equation counter: was 7, becomes 12.
%% =================================================================

\documentclass[11pt]{article}

\usepackage[margin=1in]{geometry}
\usepackage{amsmath, amssymb}
\usepackage{mathrsfs}
\usepackage{bm}
\usepackage{xcolor}
\usepackage{fancyhdr}

\pagestyle{fancy}
\fancyhf{}
\lhead{\small Published as a conference paper at ICLR 2015}
\rhead{}
\cfoot{\small 9 $\to$ \arabic{page}}
\renewcommand{\headrulewidth}{0.4pt}

\setcounter{page}{1}
\setcounter{equation}{7}

\begin{document}

% ---------- Algorithm 2 box ----------
\noindent
\fbox{\parbox{0.95\textwidth}{\small
\textbf{Algorithm 2:} \emph{AdaMax}, a variant of Adam based on
the infinity norm. See section~7.1 for details. Good default
settings for the tested machine learning problems are
$\alpha = 0.002$, $\beta_1 = 0.9$ and $\beta_2 = 0.999$. With
$\beta_1^t$ we denote $\beta_1$ to the power $t$. Here,
$(\alpha/(1 - \beta_1^t))$ is the learning rate with the
bias-correction term for the first moment. All operations on
vectors are element-wise.\\[0.6em]
\textbf{Require:} $\alpha$: Stepsize\\
\textbf{Require:} $\beta_1, \beta_2 \in [0, 1)$: Exponential decay rates\\
\textbf{Require:} $f(\theta)$: Stochastic objective function with parameters $\theta$\\
\textbf{Require:} $\theta_0$: Initial parameter vector\\[0.2em]
\hspace*{1em}$m_0 \leftarrow 0$ \hfill (Initialize $1^{\text{st}}$ moment vector)\\
\hspace*{1em}$u_0 \leftarrow 0$ \hfill (Initialize the exponentially weighted infinity norm)\\
\hspace*{1em}$t \leftarrow 0$ \hfill (Initialize timestep)\\
\hspace*{1em}\textbf{while} $\theta_t$ not converged \textbf{do}\\
\hspace*{2em}$t \leftarrow t + 1$\\
\hspace*{2em}$g_t \leftarrow \nabla_\theta f_t(\theta_{t-1})$ \hfill (Get gradients w.r.t.\ stochastic objective at timestep $t$)\\
\hspace*{2em}$m_t \leftarrow \beta_1 \cdot m_{t-1} + (1 - \beta_1) \cdot g_t$ \hfill (Update biased first moment estimate)\\
\hspace*{2em}$u_t \leftarrow \max(\beta_2 \cdot u_{t-1}, |g_t|)$ \hfill (Update the exponentially weighted infinity norm)\\
\hspace*{2em}$\theta_t \leftarrow \theta_{t-1} - (\alpha/(1 - \beta_1^t)) \cdot m_t / u_t$ \hfill (Update parameters)\\
\hspace*{1em}\textbf{end while}\\
\hspace*{1em}\textbf{return} $\theta_t$ (Resulting parameters)
}}

\vspace{1em}

% ---------- Derivation ----------
\noindent
Note that the decay term is here equivalently parameterised as
$\beta_2^p$ instead of $\beta_2$. Now let $p \to \infty$, and
define $u_t = \lim_{p \to \infty}(v_t)^{1/p}$, then:
%
\begin{align}
u_t &= \lim_{p \to \infty}(v_t)^{1/p}
     = \lim_{p \to \infty}\!\left((1 - \beta_2^p) \sum_{i=1}^{t} \beta_2^{p(t-i)} \cdot |g_i|^p\right)^{\!1/p} \\
    &= \lim_{p \to \infty}(1 - \beta_2^p)^{1/p}
       \left(\sum_{i=1}^{t} \beta_2^{p(t-i)} \cdot |g_i|^p\right)^{\!1/p} \\
    &= \lim_{p \to \infty}\!\left(\sum_{i=1}^{t} \left(\beta_2^{(t-i)} \cdot |g_i|\right)^{p}\right)^{\!1/p} \\
    &= \max\!\left(\beta_2^{t-1}|g_1|,\, \beta_2^{t-2}|g_2|,\, \ldots,\, \beta_2 |g_{t-1}|,\, |g_t|\right)
\end{align}
%
Which corresponds to the remarkably simple recursive formula:
%
\begin{equation}
u_t = \max(\beta_2 \cdot u_{t-1},\, |g_t|)
\end{equation}
%
with initial value $u_0 = 0$. Note that, conveniently enough, we
don't need to correct for initialization bias in this case. Also
note that the magnitude of parameter updates has a simpler bound
with AdaMax than Adam, namely: $|\Delta_t| \leq \alpha$.

\vspace{0.5em}

\subsection*{7.2 \textsc{Temporal averaging}}

Since the last iterate is noisy due to stochastic approximation,
better generalization performance is often achieved by
averaging. Previously in Moulines \& Bach (2011), Polyak-Ruppert
averaging (Polyak \& Juditsky, 1992; Ruppert, 1988) has been
shown to improve the convergence of standard SGD, where
$\bar{\theta}_t = \frac{1}{t} \sum_{k=1}^{n} \theta_k$.
Alternatively, an exponential moving average over the parameters
can be used, giving higher weight to more recent parameter
values. This can be trivially implemented by adding one line to
the inner loop of algorithms~1 and~2:
$\bar{\theta}_t \leftarrow \beta_2 \cdot \bar{\theta}_{t-1} + (1 - \beta_2) \theta_t$,
with $\bar{\theta}_0 = 0$. Initalization bias can again be
corrected by the estimator
$\widehat{\theta}_t = \bar{\theta}_t / (1 - \beta_2^t)$.

\vspace{0.5em}

\section*{8 \textsc{Conclusion}}

We have introduced a simple and computationally efficient
algorithm for gradient-based optimization of stochastic objective
functions. Our method is aimed towards machine learning problems
with

% [page ends here — continues on page 10]

\end{document}
